\appendix
\section{Benchmark generation cost}
\label{app:cost}

Table~\ref{tab:cost} recomputes elapsed time from the frozen training-tier
records, with 48 draws per geometry and 19 concurrent single-threaded solver
jobs. The recorded time covers imperfection construction, solver execution and
output handling. It is elapsed time under this concurrency, not measured CPU
time or an isolated-solver benchmark. The Batdorf parameter is
$Z=(L^2/Rt)\sqrt{1-\nu^2}$.

\begin{table}[tbp]
\centering
\caption{Frozen training-tier generation timings. No cached zero-time records
enter these summaries. Geometry is given as $(R/t,L/R)$.}
\label{tab:cost}
\begin{tabular}{lrrrr}
\toprule
geometry & $Z$ & nodes & mean [s] & range [s] \\
\midrule
% Generated from the frozen, hashed training-tier draws.
$110$, $0.8$ & 67 & 5311 & 1043 & 1023--1054 \\
$140$, $0.8$ & 85 & 6450 & 1319 & 1290--1342 \\
$170$, $0.8$ & 104 & 7952 & 1791 & 1757--1816 \\
$110$, $1$ & 105 & 6328 & 1362 & 1353--1385 \\
$200$, $0.8$ & 122 & 9610 & 2269 & 2230--2301 \\
$140$, $1$ & 134 & 8385 & 2028 & 1958--2115 \\
$170$, $1$ & 162 & 10082 & 2516 & 2489--2542 \\
$110$, $1.3$ & 177 & 8362 & 2079 & 2064--2097 \\
$200$, $1$ & 191 & 11935 & 3078 & 3020--3112 \\
$140$, $1.3$ & 226 & 10707 & 2924 & 2884--3021 \\
$170$, $1.3$ & 274 & 12638 & 3560 & 3531--3606 \\
$200$, $1.3$ & 322 & 15190 & 4613 & 4524--4793 \\

\bottomrule
\end{tabular}
\end{table}

An ordinary least-squares fit of log mean elapsed time against log $Z$ gives
exponent $\CostExponent$ over these twelve geometries. Node count and geometry
co-vary with $Z$ in this design; the fit neither isolates their effects nor
predicts a new mesh, solver protocol or hardware platform without validation.
The within-geometry ranges describe the measured seeds and do not imply
deterministic runtimes.

\section{Implementation configuration and reproducibility}
\label{app:config}

The checked-in SAOI profiles are later experiments, not the original sweep
reported in Section~\ref{sec:results:saoi}. Both current arm-1 profiles use two
coarse levels (1000 and 100 nodes), a $4,6,8,6,4$ message-passing schedule,
hidden width 128, four cached hierarchy variants and loss weights
$(0.1,0.1,1)$. They differ in more than stochasticity placement: the current
MGN-V arm uses 1000 training epochs and input-noise scale 0.01, while the
current cHI-MGNflow arm uses 2000 epochs and input-noise scale zero.
These values describe executable profiles, not recovered historical settings.
Architecture options and their objectives are specified in
Section~\ref{sec:method}.

The manuscript's buckling measurements are generated by
\texttt{paper/audit\_data.py} from a frozen JSON snapshot and SHA-256-identified
draw archives. This recomputes mode labels, spectra, radial displacements,
drift counts and timing tables; it does not rerun the nonlinear solves or
recover deleted force/energy histories. Plot scripts consume the same frozen
record list. The source/PDF checker verifies included files, references,
citations, figure existence and a build manifest; scientific evidence remains
bounded by the audits described in the text.

For the corrected taper runs, the archive records the solver settings,
including nodal-thickness schema, to prevent reuse of the original
uniform-thickness solves. Export requires explicit plan and tier selection,
checks the declared boundary displacement and field arrays, and records all
rejections. Each model's loader must still be configured for the documented
static information boundary. A validation/test split of draws from a known
training geometry evaluates new realisations; a held-out geometry tier
evaluates geometric transfer. These are different tasks.
